% !TEX program = xelatex
% 上行声明本文件用 XeLaTeX 编译。Overleaf 不读取该声明，按下方第 3 步设置。
% ASu-skills LaTeX 简历
%
% 供 Overleaf 等环境继续编辑排版。
%
% ── Overleaf 使用步骤（无需本地安装任何 TeX 环境）─────────────────
%   1. New Project → Blank Project，填写 Project name → Create；
%   2. 上传本 .tex 文件；若简历使用证件照或校标，把图片一并上传到同一项目
%      （与 .tex 同级，文件名保持一致）；
%   3. 左上角 File → Settings：Main document 选本文件，
%      Compiler 选 XeLaTeX（中文必需）→ Recompile。

\documentclass[a4paper,10pt]{ctexart}

\usepackage[top=1.5cm,bottom=1.5cm,left=1.6cm,right=1.6cm]{geometry}
\usepackage{enumitem}
\usepackage{xcolor}
\usepackage{titlesec}
\usepackage[hidelinks]{hyperref}
\usepackage{graphicx}

% 简历采用黑白配色，分区标题与分隔线统一为黑色
\definecolor{asublue}{HTML}{000000}

\titleformat{\section}
  {\normalfont\large\bfseries\color{asublue}}{}{0pt}{}
  [{\color{asublue}\titlerule[1.2pt]}]
\titlespacing*{\section}{0pt}{10pt}{5pt}

\setlist[itemize]{leftmargin=1.3em,itemsep=1pt,parsep=0pt,topsep=2pt,partopsep=0pt}

\setlength{\parindent}{0pt}
\pagestyle{empty}

\newlength{\asuphotowidth}
\newlength{\asuheadtextwidth}

\newcommand{\asuphoto}[1]{%
  \IfFileExists{#1}%
    {\includegraphics[width=\dimexpr\asuphotowidth-4mm\relax]{#1}}%
    {\fbox{\parbox[c][28mm][c]{\dimexpr\asuphotowidth-4mm-2\fboxsep-2\fboxrule\relax}%
       {\centering\footnotesize 未找到图片\\\texttt{#1}\\请一并上传}}}%
}

\newcommand{\asuname}[1]{{\LARGE\bfseries #1}\par\vspace{2pt}}
\newcommand{\asumeta}[1]{{\small #1}\par}
\newcommand{\asuentry}[2]{\textbf{#1}\hfill{\small #2}\par}

\begin{document}

\setlength{\asuphotowidth}{0pt}
\setlength{\asuheadtextwidth}{\dimexpr\linewidth-\asuphotowidth\relax}

\noindent
\begin{minipage}[c]{\asuheadtextwidth}
\asuname{雷明泽}
\asumeta{13892553456 \textbar{} 2454316259@qq.com}
\asumeta{2028 届硕士 · 大模型应用 / 后端开发 · 可实习 6 个月}
\end{minipage}%


\section*{教育经历}
\asuentry{北京信息科技大学 · 计算机技术（硕士）}{2025-2028}
\asuentry{陕西理工大学 · 计算机科学与技术（本科）}{2020-2024}

\section*{实习与工作经历}
\asuentry{科大讯飞股份有限公司 · 实习}{2026.04-至今}
\begin{itemize}
  \item \textbf{多模态数据链路：}基于 Data-Juicer 与脚本开发 10+ 类数据处理算子，完成协议适配、流水线编排、压测与异常排查，累计支撑千万级文本、图像、语音样本的过滤与质量评估，交付语义增强、思维链合成、问答生成和标准化导出能力。
  \item \textbf{质量治理与自动质检：}建设合成与回流数据质检链路，覆盖格式异常、重复、语义漂移和问答不匹配；开发 10+ 项清洗、校验与打分规则，实现百万级数据批量检测。
  \item \textbf{轻量模型部署与评测：}完成 10+ 个轻量模型及数据合成工具的本地部署与推理联调，覆盖情感语音合成、方言克隆、意图识别和情感分析；基于准确率、F1、生成成功率、语音自然度与相似度进行评测，为模型接入方案提供依据。
\end{itemize}

\section*{项目经历}
\textbf{仿 TikTok 的短视频平台开发}\par
面向内容社区场景，实现账号、视频、关注、点赞、评论、Feed、Top30 热榜和内容审核等核心功能。\par
\begin{itemize}
  \item \textbf{技术栈：}Go（Gin/GORM）、MySQL、Redis、RabbitMQ、WebSocket、vLLM。
  \item \textbf{推拉结合 Feed：}普通作者写扩散至粉丝 Inbox，高粉作者读时拉取并合并去重，降低写扩散压力。
  \item \textbf{热榜与异步投递：}基于 Redis ZSET 排序和快照 MD5 判变，经 Pub/Sub 通知 API 实例并通过 WebSocket 推送；点赞、评论经 RabbitMQ Worker 异步落库，视频发布使用事务 Outbox 规避 DB/MQ 双写漏发。
  \item \textbf{高频读缓存：}视频、Feed、计数和关注列表使用 Redis；Pipeline 批量读取点赞状态，singleflight 合并高粉作者回源，缓解缓存击穿。
  \item \textbf{多模态内容审核：}异步视频审核结合挑帧、OCR 和语音转写，识别切镜、插帧、半透明叠加与画中画；合成插帧测试中硬插帧拦截率由 27\% 提升至 93\%，初筛耗时由 1.1 s 降至 170 ms，约减少 53\% 的 9B 调用，并由 Qwen3.5-9B 输出结构化结论。
\end{itemize}



\section*{专业技能}
\begin{itemize}
  \item \textbf{Agent 与 RAG：}熟悉 LangChain、LangGraph，理解 ReAct、Tool Calling、Function Calling、MCP 和 Agent Harness；了解 RAG 设计与 Milvus。
  \item \textbf{编程与后端：}熟练掌握 Python、Go，能够使用 FastAPI、Gin、GORM 开发 LLM 应用及后端接口；熟练使用 MySQL、Redis、RabbitMQ，熟悉 vLLM 部署与推理。
  \item \textbf{工程化：}熟练使用 Docker、Docker Compose 和 Git；使用 Claude Code、Codex 等 AI 研发工具，能够完成需求拆解、实现、测试与部署排查。
\end{itemize}

\section*{获得奖项}
\begin{itemize}
  \item \textbf{蓝桥杯：}C++ B 组省一等奖。
\end{itemize}



\end{document}
